% !TeX root = ../navier-stokes-manuscript.tex
% ============================================================
% 1.1 The Regularity Problem
% ============================================================
\subsection{The Regularity Problem}\label{sub:TheRegularityProblem}

The three-dimensional Navier--Stokes regularity problem asks whether a smooth fluid can reach a finite time beyond which its own equations cannot continue it smoothly.
Classical local theory already starts the motion and determines it uniquely for a short time.
The unresolved part begins with the possibility that this smooth construction could have a last finite time.

This paper treats the unforced whole-space problem on \(\mathbb R^3\).
Fix a viscosity \(\nu>0\), choose smooth divergence-free initial velocity \(u_0\) with rapid decay and finite energy, and let velocity \(u(x,t)\) and pressure \(p(x,t)\) obey
\[
  \partial_tu+(u\cdot\nabla)u+\nabla p
  =
  \nu\Delta u,
  \qquad
  \nabla\cdot u=0,
  \qquad
  u(x,0)=u_0(x).
\]
Let \([0,T_*)\) be the maximal interval on which the classical solution generated by these fixed choices exists.
The existence-and-smoothness problem is exactly the claim
\[
  T_*=\infty
\]
for every admissible \(u_0\) \parencite{Fefferman2000}.

\divider{Smoothness}

\begin{definition}[Smoothness]\label{def:Smoothness}
A velocity--pressure pair \((u,p)\) is smooth through a finite time \(T>0\) when
\[
  u\in C^\infty\!\left(\mathbb R^3\times[0,T];\mathbb R^3\right),
  \qquad
  p\in C^\infty\!\left(\mathbb R^3\times[0,T]\right).
\]
Equivalently, every mixed spatial and temporal derivative of every finite order exists and is continuous through \(t=T\).
The pair is smooth for all time when this holds through every finite \(T>0\).
\end{definition}

Smoothness asks the whole field to arrive at the time slice with all of its finite derivative orders intact.
The velocity itself can approach a finite value while a gradient, an acceleration, or a higher response still fails to arrive continuously.
That is why the problem cannot be settled by watching the size of \(u\) alone.

\divider{Regularity}

\begin{definition}[Regularity]\label{def:Regularity}
Regularity tells us how much of the smooth field has actually been controlled.
The classical classes \(C^k\) require continuous mixed derivatives through total order \(k\), while Sobolev spaces measure specified spatial derivatives in integral norms.
Smoothness is the case in which every finite classical order is available.
\end{definition}

These levels matter because the equation can turn sufficiently strong finite-order control into higher control.
A bound that reaches the right Sobolev level controls the coefficients in the differentiated equation, lets higher rows be propagated from the smooth datum, and eventually returns the full smooth field.
The issue is finding such a bound from the data and the equation before a candidate endpoint is reached.

\divider{Continuation and Maximal Time}

\begin{definition}[Continuation]\label{def:Continuation}
Suppose \((u,p)\) solves the Navier--Stokes equations on \(\mathbb R^3\times[0,T)\).
A continuation past \(T\) is a pair \((\widetilde u,\widetilde p)\) solving the same equations with the same viscosity on \(\mathbb R^3\times[0,T+\varepsilon)\) for some \(\varepsilon>0\), agreeing with \((u,p)\) before \(T\), and belonging to the required regularity class on the larger interval.
\end{definition}

At any regular time \(t_0<T_*\), the complete slice \(u(\cdot,t_0)\) can be used as initial data for local theory again.
If the size of that slice in a continuation-grade class stays uniformly controlled as \(t_0\uparrow T_*\), the restarted lifespan does not collapse and a sufficiently late restart crosses \(T_*\).
Maximality can survive only if every available restart of this kind is defeated near the endpoint.
The problem is to decide whether the coupled velocity, pressure, incompressibility, transport, and viscosity can ever produce that defeat from smooth finite-energy data.
